\documentclass[10pt]{article}
\usepackage[T1]{fontenc}
\usepackage[utf8]{inputenc}
\usepackage{lmodern}
\usepackage[english]{babel}
\usepackage{amsmath,amssymb}
\usepackage[a4paper,margin=2.2cm]{geometry}
\usepackage{microtype}
\usepackage[hidelinks]{hyperref}
\hypersetup{
  hypertexnames=false,
  pdftitle={On the Invariant Theory of Forms in n Variables -- Section 5},
  pdfauthor={Emmy Noether},
  pdfsubject={Source-audited English working translation of Section 5, printed pages 134--137}
}
\newcommand{\pair}[2]{(#1\mid #2)}
\newcommand{\eps}{\varepsilon}
\setlength{\parindent}{1.25em}
\setlength{\emergencystretch}{2em}
\raggedbottom

\begin{document}

\begin{center}
{\Large\bfseries On the Invariant Theory of Forms in $n$ Variables}\par
\vspace{0.5em}
By Miss \emph{Emmy Noether} in Erlangen.\par
\medskip
\rule{0.22\linewidth}{0.4pt}\par
\vspace{0.7em}
\emph{Journal f\"ur die reine und angewandte Mathematik} 139 (1911), pp. 118--154\par
\smallskip
{\small\itshape Source-audited working draft containing \S5 only; original printed pp. 134--137.}
\end{center}

\section*{\S\ 5. The Decomposition Identities in Explicit Form.}

If, in case 4) of (25), we specialize $\tau$ to $\sigma+\rho$, we obtain
\begin{equation}
 \pair{q_\rho A^\sigma}{y^{(1)}\cdots y^{(\rho+\sigma)}}
 =\sum_{i_\rho=i_{\rho-1}+1}^{\tau}
 \eps_{i_\rho}\pair{q_\rho}{y^{(i_1)}\cdots y^{(i_\rho)}}
 \pair{A^\sigma}{y^{(i_{\rho+1})}\cdots y^{(i_{\rho+\sigma})}},
\tag{32}
\end{equation}
where $\eps_{i_\rho}=(-1)^{\delta_{i_\rho}}$ and
$\delta_{i_\rho}=\rho(\rho+1)/2+i_1+\cdots+i_\rho$. This is a more general
form of the product theorem for matrices than (16) and (17), and follows by
Laplace expansion of the determinant of the product matrix
\[
 \sum\pm(v^{(1)}y^{(1)})\cdots(v^{(\rho)}y^{(\rho)})
 (a^{(1)}y^{(\rho+1)})\cdots(a^{(\sigma)}y^{(\rho+\sigma)}).
\]
Thus the more general product theorem is composed from identities (20) and
(21), respectively, and this explains the extraction of the special forms
(16) and (17).

Relation (32) now provides the means for representing the decomposition
identities explicitly. To this end we regard the forms occurring in (31) as
ground forms; that is, we perform substitution (11), which changes nothing in
the explicit representation in the rows $A,B$, since, by (8), the newly
introduced rows $R$ can be expressed in terms of the rows $A,B$ themselves
and not in terms of their component rows $a^{(1)}a^{(2)}\ldots$,
$b^{(1)}b^{(2)}\ldots$. Thus let
\begin{equation}
 \pair{q_{\rho-\alpha}B^{n-\rho-\sigma+\tau}}{A_{n-\sigma}p_{\tau-\alpha}}
 =\bigl(R_{\rho-\alpha}p_{n-\rho+\alpha}\bigr)
  \pair{R^{\tau-\alpha}}{p_{\tau-\alpha}}.
\tag{33}
\end{equation}
Then the individual sum of (25) corresponding to the index $\alpha$ becomes
\begin{equation}
\begin{aligned}
&\sum \eps_{i_\alpha}
 \bigl(R_{\rho-\alpha}p_{n-\rho}y^{(i_1)}\cdots y^{(i_\alpha)}\bigr)
 \pair{R^{\tau-\alpha}}{y^{(i_{\alpha+1})}\cdots y^{(i_\tau)}}\\
&\quad=\sum \eps_{i_\alpha}
 \bigl([R_{\rho-\alpha}p_{n-\rho}]^\alpha
 y^{(i_1)}\cdots y^{(i_\alpha)}\bigr)
 \pair{R^{\tau-\alpha}}{y^{(i_{\alpha+1})}\cdots y^{(i_\tau)}}\\
&\quad=\eps\pair{[R_{\rho-\alpha}p_{n-\rho}]^\alpha
 R^{\tau-\alpha}}{p_\tau}
 =\eps\pair{R^{\tau-\alpha}q_{n-\tau}}
 {R_{\rho-\alpha}p_{n-\rho}}.
\end{aligned}
\tag{34}
\end{equation}
An explicit final expression has therefore indeed been obtained, and the
symmetric form in which $q_\rho$ and $p_\tau$ occur shows that (30) leads to
the same final expression, which is consequently independent of the original
decomposition. A difference between (25) and (30) exists only as long as the
rows $y$ and $v$ have independent significance, so that the decomposition
identity passes, according to our definition, into a decomposable identity.
To obtain identities with more symbol rows, one need only, by \S\ 3 (end),
resolve a row $q_\rho$ into $q_{\rho_1}C^{\sigma_1}$, or $p_\tau$ into
$p_{\tau_1}C_{\tau-\tau_1}$. The resulting expressions
\begin{equation}
 \pair{q_{\rho_1}C^{\sigma_1}}
 {[R^{\tau-\alpha}q_{n-\tau}]_\lambda
  R_{\rho_1+\sigma_1-\alpha-\lambda}},
 \qquad
 \pair{[R_{\rho-\alpha}p_{n-\rho}]^\lambda R^{\tau-\alpha}}
 {p_{\tau_1}C_{\tau-\tau_1}}
\tag{35}
\end{equation}
are exactly of the type occurring with two symbol rows, and therefore again
admit explicit representation after a substitution corresponding to (33).
Repetition consequently yields an explicit representation for arbitrarily
many rows. It may also be noted that the first and last sums in (25) and
(30), respectively, yield the explicit representation without applying
(33), solely by formula (32), or else reduce altogether to a single term
containing all rows explicitly. The relation obtained from (32) by resolving
$q_\rho$ into $q_{\rho_1}C^{\sigma_1}$, and so forth, may be regarded as the
product theorem in its most general form.\footnote{In the original print,
the two expressions in (35) use the component index $\alpha$, with the first
containing
$[R^{\tau-\alpha}q_{n-\tau}]_\alpha
R_{\rho_1+\sigma_1-\alpha}$. The current R823 German source makes the
coordinated editorial change from $\alpha$ to $\lambda$ in both expressions
and correspondingly uses
$[R^{\tau-\alpha}q_{n-\tau}]_\lambda
R_{\rho_1+\sigma_1-\alpha-\lambda}$ in the first; R823 is followed here.}

In summary, we shall show:

\emph{Theorem III.} With the identities
\begin{equation}
 \pair{q_\rho A^\sigma}{p_\tau B_{\rho+\sigma-\tau}}
 =\sum_{\substack{0,\tau-\sigma}}^{\substack{\tau,\rho}}
 \eps\pair{R^{\tau-\alpha}q_{n-\tau}}
 {R_{\rho-\alpha}p_{n-\rho}},
\tag{36}
\end{equation}
where the rows $R$ are determined by (33), together with those obtained from
them by resolving $q_\rho,p_\tau$ into $q_{\rho_1}C^{\sigma_1}$, and so
forth, the totality of indecomposable identities is exhausted. In this
combined-bound notation, as in the preceding section, the summation index
runs from the larger lower value to the smaller upper value.

Indeed, the resulting identities are the most general of their kind, since
the individual rows are no longer subject to any restriction with respect to
weight or number, as was still the case in (20) and (21). Nor do any further
identities among these rows exist. For, on resolving a row $p_\tau$ into
$y^{(1)}\cdots y^{(\tau)}$, the most general identities are composed from
(20), and hence, by \S\ 3 (end), from (20a); the composition discussed after
(16) is precisely the one carried out in \S\ 4. The passage to arbitrarily
many rows is supplied, as the discussion of (20a) shows, by applying the
same operation repeatedly to the expressions arising from the resolution of
$q_\rho$ into $q_{\rho_1}C^{\sigma_1}$, and so forth. It follows as well that
the direct passage from (20), instead of from (20a), yields nothing new. This
is easily verified computationally: specialize a row $q_\rho$ in (25) once,
and another time carry out the passage from (20) by the method of \S\ 4. The
same sums arise both times, and application of the general product theorem
(see above) transforms them into the identities obtained from (36). The
identities (20) and (21) correspond to the special cases $\tau=1$ and
$\rho=1$; only two individual sums then occur in each case, hence, in
accordance with the preceding remark, an explicit representation without
substitution (33), in agreement with the earlier result.

Finally, let us briefly mention two special cases of the decomposition
identity. In (31), set $\tau=\sigma$, $\rho=n-\sigma$, and denote the row
$p_\tau$ by $p'_\sigma\sim q'_{n-\sigma}$. Then
\begin{equation}
\begin{aligned}
\pair{A^\sigma}{p_\sigma}\pair{B^\sigma}{p'_\sigma}
-\pair{A^\sigma}{p'_\sigma}\pair{B^\sigma}{p_\sigma}
=\sum_{\alpha=1}^{\sigma,n-\sigma}
\Pi\pair{q_{n-\sigma-\alpha}B^\sigma}
{A_{n-\sigma}p_{\sigma-\alpha}}.
\end{aligned}
\tag{37}
\end{equation}
This is Clebsch's formula\footnote{Clebsch, loc. cit., pp. 25--32. Clebsch
shows that the individual terms on the right-hand sides depend only on the
rows $A,B$; the explicit representation would result by applying (32) to his
expressions $\Phi_h$.} for expanding a form with two cogredient rows
$p_\sigma,p'_\sigma$ in elementary covariants. The elementary covariants
belonging to a form
$\pair{A^\sigma}{p_\sigma}^m\pair{B^\sigma}{p'_\sigma}^n$ are then the
following:
\begin{equation}
 \left\{\sum_{\alpha=1}^{\sigma,n-\sigma}
 \pair{q_{n-\sigma-\alpha}B^\sigma}
 {A_{n-\sigma}p_{\sigma-\alpha}}
 \right\}^{\rho}
 \pair{A^\sigma}{p_\sigma}^{m-\rho}
 \pair{B^\sigma}{p'_\sigma}^{n-\rho},
 \qquad (\rho=0,1,\ldots,m,n).
\tag{38}
\end{equation}
Here the paired upper endpoints again mean the smaller of the two displayed
values; likewise, the final range for $\rho$ ends at the smaller of $m$ and
$n$. These covariants arise, as we shall briefly say, from the ground form by
``cogredient folding of defect $\rho$'' (cf. \S\ 2).\footnote{The original
print reads ``defect $\alpha$'' in this sentence. The current R823 German
source reads ``defect $\rho$''; R823 is followed here because $\rho$ is the
outer exponent and the immediately stated range parameter in (38). This is a
disclosed print/R823 editorial delta, not a claim that the printed reading is
mathematically erroneous.}

Second, set $\tau=\sigma-\lambda$, $\rho=n-\sigma-\lambda$, and
$B^\sigma=A^\sigma$. Then
\begin{equation}
\begin{aligned}
\pair{q_{n-\sigma-\lambda}A^\sigma}
{A_{n-\sigma}p_{\sigma-\lambda}}
&=(-1)^\lambda\pair{q_{n-\sigma-\lambda}A^\sigma}
{A_{n-\sigma}p_{\sigma-\lambda}}\\
&\quad+(-1)^{(n-\sigma)(\sigma-\lambda)}
 \sum_{\alpha=1}^{\sigma-\lambda,n-\sigma-\lambda}
 \Pi\pair{q_{n-\sigma-\lambda-\alpha}A^\sigma}
 {A_{n-\sigma}p_{\sigma-\lambda}}.
\end{aligned}
\tag{39}
\end{equation}
Thus, as noted in \S\ 2, for odd defect $\lambda$ the conditions (12)
characterizing the symbol rows are consequences of the conditions of higher
defect. For a linear form
$\pair{A^\sigma}{p_\sigma}=\pair{B^\sigma}{p_\sigma}$,\footnote{The original
print places the identity sign $\equiv$ between these two expressions; the
current R823 German source uses ordinary equality. R823 is followed here, with
the historical notation preserved in this disclosure. This is recorded as an
editorial notation delta, not as a proved mathematical source defect.} it follows from (39)
that its irreducible invariant formations that are quadratic in the
coefficients reduce to ones of even defect. This agrees with the well-known
fact that a linear form in the binary and ternary domains possesses no
invariant formations.

It should be noted that the general identities were originally derived under
the assumption that the individual rows satisfy conditions (12). Since,
however, these individual rows enter the identities only linearly, the latter
hold also for the more general rows that do not satisfy conditions (12).

\end{document}
